\documentclass{article}
\usepackage{amsmath,amssymb}
\usepackage{booktabs}
\usepackage[margin=1in]{geometry}
\title{Step 227: Bridge Impossibility on the Hodge Track}
\author{Six Birds Foundations III}
\date{2026-05-16}
\begin{document}
\maketitle

\section{Purpose}

This note tests whether the Bridge Impossibility Corollary, first derived on the RH track and verified on BSD, replicates for the Hodge conjecture.

\section{Hodge Track Records}

The Hodge cascade map records the active parent residual
\[
  \Xi_H^{\rm std}(X^4_{33},2)\neq 0
\]
as a structured standard-dictionary residual. The track's reference-audit step H-18R states the source-to-status rule:
\[
  \texttt{known\_zero}\quad\text{is licensed only by a source-stated cycle-span theorem}\quad A_{X,p}=V_{X,p}.
\]
Thus arithmetic support, variational structure, or projector architecture is not a verified Hodge source column until the cycle dictionary and projection/descent gates are supplied.

\section{Hodge Bridge Impossibility Analog}

\textbf{Claim (Hodge Bridge Impossibility, cross-track test).}
Let \(C_H\) be a carrier with a proved analog \(\Xi_{C_H}=0\) that is not equivalent to the full Hodge conjecture. Any typed bridge \(B\) that transports this closure to full Hodge makes the composite \(C_H+B\) Hodge-strength.

\section{Candidate Carriers}

\begin{center}
\begin{tabular}{p{0.23\linewidth}p{0.27\linewidth}p{0.28\linewidth}p{0.12\linewidth}}
\toprule
Carrier & Proved scope & Bridge to full Hodge & Assessment\\
\midrule
Lefschetz (1,1) & Codimension-one divisor classes & Higher codimension algebraicity & Hodge-strength\\
Surfaces & Surface classes reduce to divisors & Dimension at least three and codimension at least two & Hodge-strength\\
Abelian varieties & Special family/class results & All smooth projective varieties & Hodge-strength\\
Standard conjectures & Projectors and positivity architecture & Algebraicity of all Hodge classes & Hodge-strength\\
Tate conjecture & l-adic analog in scoped cases & Comparison plus cycle realization & Hodge-strength / conjectural\\
CDK Hodge loci & Algebraicity of Hodge loci & Fixed-variety cycle-span theorem & Hodge-strength\\
Schubert dictionaries & Grassmannian/flag known-zero cases & Arbitrary varieties & Hodge-strength\\
\bottomrule
\end{tabular}
\end{center}

\section{Literature Audit}

The audit covered Hodge's 1950 ICM formulation, Lefschetz's \((1,1)\) theorem, Grothendieck's standard conjectures, Deligne's 1982 work on Hodge cycles on abelian varieties, Cattani--Deligne--Kaplan's theorem on Hodge loci, Voisin's expository account, Lewis's survey, and Tate's l-adic analog. These sources provide scoped closures or structural evidence. None supplies a non-target-strength bridge from a non-Hodge-equivalent carrier to full Hodge.

\section{Conclusion}

The Hodge analog is verified. The framework finding upgrades to:
\[
  \texttt{Bridge Impossibility verified-on-3-track-instances}
\]
for RH, BSD, and Hodge.

\end{document}
